%%%%%%%%%%%%%%%%%%%%%%%%%%%%%%%%%%%%%%%%%
% University Assignment Title Page 
% LaTeX Template
% Version 1.0 (27/12/12)
%
% This template has been downloaded from:
% http://www.LaTeXTemplates.com
%
% Original author:
% WikiBooks (http://en.wikibooks.org/wiki/LaTeX/Title_Creation)
%
% License:
% CC BY-NC-SA 3.0 (http://creativecommons.org/licenses/by-nc-sa/3.0/)
% 
% Instructions for using this template:
% This title page is capable of being compiled as is. This is not useful for 
% including it in another document. To do this, you have two options: 
%
% 1) Copy/paste everything between \begin{document} and \end{document} 
% starting at \begin{titlepage} and paste this into another LaTeX file where you 
% want your title page.
% OR
% 2) Remove everything outside the \begin{titlepage} and \end{titlepage} and 
% move this file to the same directory as the LaTeX file you wish to add it to. 
% Then add \input{./title_page_1.tex} to your LaTeX file where you want your
% title page.
%
%%%%%%%%%%%%%%%%%%%%%%%%%%%%%%%%%%%%%%%%%
%\title{Title page with logo}
%----------------------------------------------------------------------------------------
%	PACKAGES AND OTHER DOCUMENT CONFIGURATIONS
%----------------------------------------------------------------------------------------

\documentclass[12pt]{article}
\usepackage[italian]{babel}
\usepackage[utf8x]{inputenc}
\usepackage{amsmath}
\usepackage{graphicx}
\usepackage[colorinlistoftodos]{todonotes}

\begin{document}

\begin{titlepage}

\newcommand{\HRule}{\rule{\linewidth}{0.5mm}} % Defines a new command for the horizontal lines, change thickness here

\center % Center everything on the page
 
%----------------------------------------------------------------------------------------
%	HEADING SECTIONS
%----------------------------------------------------------------------------------------

\textsc{\LARGE Università degli studi di Milano-Bicocca}\\[1cm] % Name of your university/college
\textsc{\Large Advanced Machine Learning }\\[0.3cm] % Major heading such as course name
\textsc{\large Progetto Finale}\\[0.1cm] % Minor heading such as course title

%----------------------------------------------------------------------------------------
%	TITLE SECTION
%----------------------------------------------------------------------------------------

\HRule \\[0.4cm]
{ \huge \bfseries Una rete neurale convoluzionale patch-based per la prevenzione della retinopatia diabetica}\\[0.4cm] % Title of your document
\HRule \\[1.5cm]
 
%----------------------------------------------------------------------------------------
%	AUTHOR SECTION
%----------------------------------------------------------------------------------------

\large
\emph{Autori:}\\
Lorenzo Mammana - 807391- l.mammana@campus.unimib.it \\   % Your name
Eric Nisoli - 807147- e.nisoli1@campus.unimib.it   \\[1cm] % Your name

% If you don't want a supervisor, uncomment the two lines below and remove the section above
%\Large \emph{Author:}\\
%John \textsc{Smith}\\[3cm] % Your name

%----------------------------------------------------------------------------------------
%	DATE SECTION
%----------------------------------------------------------------------------------------

{\large \today}\\[2cm] % Date, change the \today to a set date if you want to be precise

%----------------------------------------------------------------------------------------
%	LOGO SECTION
%----------------------------------------------------------------------------------------

\includegraphics{logo.png}\\[1cm] % Include a department/university logo - this will require the graphicx package
 
%----------------------------------------------------------------------------------------

\vfill % Fill the rest of the page with whitespace

\end{titlepage}


\begin{abstract}
The ABSTRACT is not a part of the body of the report itself. Rather, the abstract is a brief summary of the report contents that is often separately circulated so potential readers can decide whether to read the report. The abstract should very concisely summarize the whole report: why it was written, what was discovered or developed, and what is claimed to be the significance of the effort. The abstract does not include figures or tables, and only the most significant numerical values or results should be given.
\end{abstract}

\section{Introduzione}
La retinopatia diabetica è la causa principale di cecità tra le persone affette da qualsiasi tipologia di malattia diabetica. Se trattata precocemente è possibile ridurre la possibilità di perdere la vista fino al 90\%.
Il problema principale è che nei paesi industrializzati solamente una persona su due si sottopone a controlli regolari di prevenzione e questa percentuale si riduce ulteriormente in zone rurali, sia per la non diligenza dei paziente che, soprattutto, per l'assenza di personale in grado di sopperire a tutte le richieste.
Sviluppare un modello in grado di aiutare o ancora meglio superare le capacità di un medico nel valutare immagini retinografiche potrebbe permettera la diagnosi precoce di questa malattia anche in aree rurali salvando quindi la vista ad un numero molto elevato di persone.
Attualmente lo stato dell'arte per un compito di questo tipo è rappresentato dalle reti neurali convoluzionali (CNN), questi modelli si sono dimostrati efficienti nel classificare immagini negli ambiti più disparati, andando molto spesso a superare le capacità discriminative di esperti nel settore.
Una delle più grosse limitazioni delle CNN è rappresentato dalla capacità in termini di memoria delle moderne GPU, uno dei compromessi che di solito viene utilizzato è quello di ridurre la dimensione delle immagini utilizzate per permettere l'utilizzo di una rete con un numero molto elevato di parametri, in grado potenzialmente di operare meglio di una rete con un numero di parametri minore.
Le immagini retinografiche in nostro possesso presentano una risoluzione elevata, molti studi hanno dimostrato come in ambito medicale sia generalmente infattibile ridurre la dimensione delle immagini, in quanto questo porta alla perdita di dettagli di piccola dimensione, fondamentali per ottenere una buona discriminazione.
Per questo motivo abbiamo deciso di utilizzare un modello \textit{patch-based} di rete neurale convoluzionale; modelli di questo tipo permettono di ottenere un compromesso migliore tra capacità computazionale della rete e dimensione dell'immagine in ingresso.
L'idea è quella di sfruttare l'intera immagine retinopatica senza perdere dettagli importanti cercando di discriminare quali sono le parti interessanti dell'immagine per poter poi usarle in maniera efficace per il compito di classificazione.
The introduction should provide a clear statement of the problem posed by the project, and why the problem is of interest. It should reflect the scenario, if available. If needed, the introduction also needs to present background information so that the reader can understand the significance of the problem. A brief summary of the hypotheses and the approach your group used to solve the problem should be given, possibly also including a concise introduction to theory or concepts used later to analyze and to discuss the results.


\section{Dataset}
Il dataset è stato scaricato dalla piattaforma kaggle (mettere link).
E' stato costruito dal \textit{Aravind Eye Hospital} ubicato in India e contiene immagini di fondo oculare di più di 5000 pazienti.
Le immagini hanno dimensione differente, la maggior parte sono ad alta risoluzione (>3000x2000). Ogni immagine presenta una label così definita:
\begin{itemize}
\item 0 - assenza di retinopatia diabetica
\item 1 - retinopatia diabetica leggera
\item 2 - retinopatia diabetica moderata
\item 3 - retinopatia diabetica severa
\item 4 - retinopatia diabetica proliferativa
\end{itemize}
Come si evince facilmente, l'ordine delle label rappresenta la gravità della malattia.

In this section the available data sets must be presented. The term dataset refers to any type of information source, for example web services for geolocation fall into this category. 
In addition, all necessary data manipulation processes, such as cleaning and enrichment with external sources, must be presented and discussed.

\section{The Methodological Approach}

This is the central and most important section of the report. Its objective must be to show, with linearity and clarity, the steps that have led to the definition of a decision model. The description of the working hypotheses, confirmed or denied, can be found in this section together with the description of the subsequent refining processes of the models. Comparisons between different models (e.g. heuristics vs. optimal models) in terms of quality of solutions, their explainability and execution times are welcome. 

Do not attempt to describe all the code in the system, and do not include large pieces of code in this section, use pseudo-code where necessary. Complete source code should be provided separately (in Appendixes, as separated material or as a link to an on-line repo). Instead pick out and describe just the pieces of code which, for example:
\begin{itemize}
\item are especially critical to the operation of the system;
\item you feel might be of particular interest to the reader for some reason;
\item  illustrate a non-standard or innovative way of implementing an algorithm, data
structure, etc..
\end{itemize}

You should also mention any unforeseen problems you encountered when implementing the
system and how and to what extent you overcame them. Common problems are:
 difficulties involving existing software.


\section{Results and Evaluation}
The Results section is dedicated to presenting the actual results (i.e. measured and calculated quantities), not to discussing their meaning or interpretation. The results should be summarized using appropriate Tables and Figures (graphs or schematics). Every Figure and Table should have a legend that describes concisely what is contained or shown. Figure legends go below the figure, table legends above the table. Throughout the report, but especially in this section, pay attention to reporting numbers with an appropriate number of significant figures. 

\section{Discussion}
The discussion section aims at interpreting the results in light of the project's objectives. The most important goal of this section is to interpret the results so that the reader is informed of the insight or answers that the results provide. This section should also present an evaluation of the particular approach taken by the group. For example: Based on the results, how could the experimental procedure be improved? What additional, future work may be warranted? What recommendations can be drawn?


\section{Conclusions}
Conclusions should summarize the central points made in the Discussion section, reinforcing for the reader the value and implications of the work. If the results were not definitive, specific future work that may be needed can be (briefly) described. The conclusions should never contain ``surprises''. Therefore, any conclusions should be based on observations and data already discussed. It is considered extremely bad form to introduce new data in the conclusions.

\section*{References}

The references section should contain complete citations following standard form.  The references should be numbered and listed in the order they were cited in the body of the report. In the text of the report, a particular reference can be cited by using a numerical number in brackets as \cite{Lee2015} that corresponds to its number in the reference list. \LaTeX provides several styles to format the references

\bibliographystyle{IEEEtran}
\bibliography{references.bib}

\end{document}